\section*{Overview}

Digit DP is the technique for counting integers in a large range when the property depends on the digits themselves.
Instead of iterating through every number up to \(10^{18}\), I build the answer one digit at a time while tracking
just enough state to know whether a partial prefix can still lead to a valid number.

\section*{When to Use It}

Use it when:

\begin{itemize}
  \item the range endpoint is huge, so direct iteration is impossible,
  \item validity depends on decimal digits, digit counts, adjacency, sums, or remainders,
  \item the natural way to answer \([L, R]\) is \(\texttt{solve}(R) - \texttt{solve}(L - 1)\).
\end{itemize}

If the property does not care about digits specifically, a normal combinatorial DP or prefix technique is usually
cleaner.

\section*{Core Idea}

Write an upper bound \(N\) as a string of digits. Process it from left to right.

The standard state usually contains:

\begin{itemize}
  \item \textbf{position}: which digit are we filling now,
  \item \textbf{tight}: have we matched the prefix of \(N\) exactly so far,
  \item \textbf{started}: have we placed a non-leading-zero digit yet,
  \item problem-specific state such as digit sum, remainder, previous digit, or count information.
\end{itemize}

The \texttt{tight} flag is the reason the DP respects the upper bound. If it is still true, the current digit cannot
exceed the corresponding digit of \(N\). Once it becomes false, the rest of the positions can use any digit.

\section*{Operations}

\begin{itemize}
  \item define \(\texttt{solve}(N)\) on one upper bound,
  \item memoize only the states with \texttt{tight = false},
  \item loop over the next digit and update the problem-specific state,
  \item answer a range by subtraction.
\end{itemize}

The reference code counts numbers whose digit sum is exactly a given target. That example is small enough to read,
but the same skeleton extends to many other properties.

\section*{Correctness Intuition}

Every integer in \([0, N]\) corresponds to exactly one root-to-leaf path in the state graph: the path that writes its
decimal digits from left to right. The \texttt{tight} flag guarantees that invalid prefixes above \(N\) are never
generated, and the problem-specific state records all information needed to decide validity at the end.

Because two states with the same \((position, problem\_state, started)\) and \(\texttt{tight = false}\) have exactly
the same future options, memoizing them is valid.

\section*{Complexity Analysis}

If the problem-specific state size is \(S\) and the decimal alphabet size is \(10\), then the complexity is roughly
\(O(\text{digits} \cdot S \cdot 10)\).

For the example note code:

\begin{itemize}
  \item states: position, sum so far, started,
  \item transitions per state: 10 digits,
  \item memory: \(O(\text{digits} \cdot \text{target sum})\).
\end{itemize}

\section*{Implementation}

The reference implementation exposes:

\begin{itemize}
  \item \texttt{count\_exact\_sum(limit, wanted\_sum)}
  \item \texttt{count\_in\_range(left, right, wanted\_sum)}
\end{itemize}

It treats leading zeros carefully so that the digit sum of \(0\) is still defined sensibly. In real problems I often
replace the sum state with a remainder, a previous digit, a bitmask, or some compressed automaton state.

\section*{Common Pitfalls}

\begin{itemize}
  \item Forgetting the \texttt{started} flag and accidentally counting leading zeros as real digits.
  \item Memoizing tight states, which depend on the exact upper bound prefix.
  \item Returning the wrong base case for the number \(0\).
  \item Letting the problem-specific state grow too large and turning the DP into a memory sink.
  \item Writing the transitions correctly for \(\texttt{solve}(N)\) but forgetting to subtract \(\texttt{solve}(L - 1)\).
\end{itemize}

\section*{Practice Problems}

\begin{itemize}
  \item Count numbers with a fixed digit sum or digit remainder condition.
  \item Count numbers with no adjacent equal digits.
  \item Range counting where the digit property can be represented by a small automaton or previous-digit state.
\end{itemize}

\section*{References}

\begin{itemize}
  \item \href{https://usaco.guide/gold/digit-dp?lang=cpp}{USACO Guide: Digit DP}.
  \item \href{https://usaco.guide/PAPS.pdf}{Principles of Algorithmic Problem Solving}.
  \item \href{https://codeforces.com/catalog}{Codeforces Catalog}.
\end{itemize}
